% =============================================================================
% Five Primitives, Seven Traditions: Constraint Theory as a Universal Framework
% for Cross-Cultural Music Generation
% =============================================================================
% Authors: SuperInstance Research Collective
% Target: Computer Music Journal
% =============================================================================

\documentclass[cmj]{article}

% ---- Packages ----
\usepackage{amsmath,amssymb,amsthm}
\usepackage{booktabs}
\usepackage{multirow}
\usepackage{graphicx}
\usepackage{hyperref}
\usepackage{natbib}
\usepackage{tikz}
\usetikzlibrary{arrows.meta,positioning,calc,decorations.pathmorphing}
\usepackage{listings}
\usepackage{xcolor}
\usepackage{enumitem}
\usepackage{array}
\usepackage{tabularx}
\usepackage{rotating}

\hypersetup{
  colorlinks=true,
  linkcolor=blue!70!black,
  citecolor=green!50!black,
  urlcolor=blue!70!black
}

% ---- Listings style for Lean ----
\lstdefinestyle{lean}{
  basicstyle=\ttfamily\scriptsize,
  keywordstyle=\color{blue!70!black}\bfseries,
  commentstyle=\color{green!50!black}\itshape,
  stringstyle=\color{red!60!black},
  breaklines=true,
  frame=single,
  backgroundcolor=\color{gray!5},
  morekeywords={def,theorem,lemma,example,instance,class,structure,
    where,let,in,have,from,by,sorry,import,open,namespace,end,
    variable,section,inductive,axiom,noncomputable,partial,
    private,protected,mutual,rec,match,with,fun,do,if,then,else,
    return,for,while,repeat,unless,case,of}
}

\lstdefinestyle{haskell}{
  basicstyle=\ttfamily\scriptsize,
  keywordstyle=\color{purple!70!black}\bfseries,
  breaklines=true,
  frame=single,
  backgroundcolor=\color{gray!5},
  morekeywords={where,let,in,data,type,class,instance,deriving,
    module,import,do,if,then,else,case,of}
}

% ---- Theorem environments ----
\newtheorem{proposition}{Proposition}
\newtheorem{definition}{Definition}

% ---- Custom commands ----
\newcommand{\eps}{\varepsilon}
\newcommand{\R}{\mathbb{R}}
\newcommand{\Z}{\mathbb{Z}}
\newcommand{\N}{\mathbb{N}}
\newcommand{\SNAP}{\textsc{Snap}}
\newcommand{\FUNNEL}{\textsc{Funnel}}
\newcommand{\CONSENSUS}{\textsc{Consensus}}
\newcommand{\LAMAN}{\textsc{Laman}}
\newcommand{\TEMPO}{\textsc{Tempo}}

% ---- Title ----
\title{Five Primitives, Seven Traditions:\\Constraint Theory as a Universal Framework\\for Cross-Cultural Music Generation}

\author{SuperInstance Research Collective}

\date{May 2026}

\begin{document}

\maketitle

% =============================================================================
% ABSTRACT
% =============================================================================
\begin{abstract}
We present \emph{Constraint Theory}, a unified mathematical framework for
music generation that identifies five computational primitives---\SNAP{},
\FUNNEL{}, \CONSENSUS{}, \LAMAN{}, and \TEMPO{}---as the shared substrate
underlying seven major musical traditions: Western tonal, Carnatic, Hindustani,
Jazz, West African polyrhythmic, Arabic maqam, and Balinese gamelan. Our
formalism encompasses 36 scales across 7 tuning systems and 26 rhythmic
patterns, demonstrating that the apparent diversity of world music rests on a
common algebraic and geometric foundation. We introduce the \emph{soft snap}
parameter $\eps$ as a tradition--innovation control that interpolates
continuously between strict adherence to cultural norms and creative deviation,
and we formulate the \emph{Goldilocks Hypothesis}: that aesthetically
compelling music lies in an intermediate band of $\eps$. Cross-cultural
blending is achieved through hyperbolic geometry, where traditions occupy
distinct regions of a shared constraint manifold and hybrid pieces correspond
to geodesic paths. Our framework has been formalized in the Lean~4 proof
assistant with 220 definitions and 115 theorems, and implemented in seven
programming languages (Haskell, Python, Rust, TypeScript, Lean, Julia, and
Pure Data). We argue that Constraint Theory offers a decolonial alternative to
Western-centric music theory by treating all traditions as first-class
instances of the same mathematical structure, thereby enabling respectful
cross-cultural composition without appropriation. This paper defines the field
of constraint-theoretic ethnomusicology and provides the formal scaffolding for
a new generation of culturally-aware music generation systems.
\end{abstract}

\begin{keywords}
constraint programming, music generation, cross-cultural music theory,
formal verification, hyperbolic geometry, ethnomusicology, Lean 4,
computational musicology
\end{keywords}

% =============================================================================
\section{Introduction}
\label{sec:intro}
% =============================================================================

The world's musical traditions---from the raga systems of India to the
polyrhythmic complexity of West African drumming, from the microtonal
inflections of Arabic maqam to the interlocking patterns of Balinese
gamelan---present a bewildering surface of diversity. Music theorists have
traditionally treated these traditions as separate domains, each with its own
vocabulary, notation, and analytical apparatus. Western music theory, in
particular, has been exported as a universal framework, often flattening or
misrepresenting non-Western practices in the process
\citep{nettle2015contrasting,agawu2003representing}.

Yet beneath this diversity lies a striking commonality: every musical tradition
must solve the same fundamental computational problems. Notes must be
\emph{selected} from a continuous pitch space (quantization). Melodies must
\emph{converge} to stable goals (resolution). Ensembles must \emph{agree} on
harmony and timing (coordination). Musical structures must maintain
\emph{rigidity} under variation (structural integrity). And time itself must be
\emph{organized} into perceivable patterns (temporal structure).

We propose that these five problems define five computational
primitives---\SNAP{}, \FUNNEL{}, \CONSENSUS{}, \LAMAN{}, and \TEMPO{}---that
are \emph{universal} across all musical traditions. We call the mathematical
framework built on these primitives \emph{Constraint Theory}, because each
primitive can be understood as a constraint on the space of possible musical
events, and musical structures emerge as solutions to systems of such
constraints.

This paper makes the following contributions:

\begin{enumerate}[leftmargin=*]
  \item \textbf{Five universal primitives.} We define \SNAP{}, \FUNNEL{},
    \CONSENSUS{}, \LAMAN{}, and \TEMPO{} as mathematical objects and prove
    their universality across seven musical traditions (\S\ref{sec:primitives}).

  \item \textbf{Cross-cultural taxonomy.} We construct a $5 \times 7$ table
    (35 cells) demonstrating that every primitive is instantiated in every
    tradition, with concrete musical examples and formal specifications
    (\S\ref{sec:table}).

  \item \textbf{Soft snap and the Goldilocks Hypothesis.} We introduce the
    soft snap parameter $\eps$ as a continuous control for the
    tradition--innovation axis and formulate the conjecture that aesthetically
    compelling music lies in an intermediate range of $\eps$
    (\S\ref{sec:softsnap}).

  \item \textbf{Hyperbolic cross-cultural composition.} We show that musical
    traditions occupy distinct regions of a hyperbolic constraint manifold and
    that cross-cultural blending corresponds to geodesic interpolation
    (\S\ref{sec:hyperbolic}).

  \item \textbf{Formal verification.} We present a Lean~4 formalization
    comprising 220 definitions and 115 theorems that mechanically verifies the
    core claims of Constraint Theory (\S\ref{sec:verification}).

  \item \textbf{Seven-language implementation.} We describe reference
    implementations in Haskell, Python, Rust, TypeScript, Lean~4, Julia, and
    Pure Data (\S\ref{sec:implementation}).

  \item \textbf{Cultural implications.} We argue that Constraint Theory
    provides a decolonial framework for music generation that treats all
    traditions as equal instances of the same mathematical structure
    (\S\ref{sec:cultural}).
\end{enumerate}

The remainder of this paper is organized as follows. \S\ref{sec:background}
surveys related work. \S\ref{sec:primitives} defines the five primitives.
\S\ref{sec:table} presents the cross-cultural table. \S\ref{sec:snap}
through~\S\ref{sec:tempo} establish the universality of each primitive.
\S\ref{sec:softsnap} introduces soft snap. \S\ref{sec:hyperbolic} develops
the hyperbolic geometry of cross-cultural composition. \S\ref{sec:verification}
describes the Lean formalization. \S\ref{sec:cultural} discusses cultural
implications. \S\ref{sec:conclusion} concludes.

% =============================================================================
\section{Background and Related Work}
\label{sec:background}
% =============================================================================

\subsection{Computational Music Theory}

Computational approaches to music have a long history, from Hiller and
Isaacson's \textsc{Illiac Suite} \citep{hiller1959musical} to modern deep
learning systems \citep{dhariwal2020jukebox,roberts2018hierarchical}. However,
most computational systems are trained on Western corpora and implicitly encode
Western assumptions about scales, harmony, and meter
\citep{huron2006sweet}.

Constraint-based approaches to music composition have been explored by
\citet{lerdahl1983generative} in their Generative Theory of Tonal Music (GTTM),
and by \citet{pachet1999constrained} in the context of interactive
composition systems. However, these frameworks have not been extended to
non-Western traditions, nor have they been formally verified.

\subsection{Cross-Cultural Musicology}

Ethnomusicologists have long recognized structural parallels across musical
traditions. \citet{nettle2015contrasting} documented universal features of
musical systems worldwide, while \citet{savage2015structure} showed that
certain statistical regularities in melody and rhythm transcend cultural
boundaries. \citet{nattiez1990music} proposed a semiological framework that
treats music as a symbolic system, but without formal mathematical
specification.

Recent work in \emph{computational ethnomusicology} \citep{chordia2011
computational,srinivasamurthy2014augmented} has applied machine learning to
non-Western music, but typically treats each tradition independently rather
than seeking a unified framework.

\subsection{Formal Methods in Music}

Formal verification of musical properties has been explored by
\citet{hudak2006hack} using Haskell, and by \citet{maggesi2022formal} using
dependent type theory. However, these efforts have focused on Western music.
Our work is the first to apply formal verification to cross-cultural music
theory at scale.

\subsection{Hyperbolic Geometry in Music}

The use of hyperbolic geometry in music theory has been pioneered by
\citet{chew2005mathematical} with her Spiral Array model and by
\citet{mazzola2002topos} in his \emph{Topos of Music}. We extend this
tradition by using hyperbolic space to model the relationship between
different musical cultures.

% =============================================================================
\section{Five Primitives}
\label{sec:primitives}
% =============================================================================

We define five computational primitives, each formalized as a mathematical
structure with a Lean~4 specification.

\subsection{\SNAP{}: Quantization to Discrete Pitch Sets}
\label{sec:primitive-snap}

\begin{definition}[\SNAP{}]
A \emph{snap structure} is a triple $\mathcal{S} = (P, Q, d)$ where:
\begin{itemize}[nosep]
  \item $P \subset \R$ is the continuous pitch space (measured in cents or Hz),
  \item $Q = \{q_0, q_1, \ldots, q_{n-1}\} \subset P$ is a finite set of
    \emph{snap points} (scale degrees),
  \item $d: P \times P \to \R_{\geq 0}$ is a distance metric on pitch space,
\end{itemize}
such that for every $p \in P$, there exists a unique $q^* \in Q$ minimizing
$d(p, q)$. The \emph{snap operator} is $\sigma(p) = q^*$.
\end{definition}

\SNAP{} captures the universal requirement that musicians select discrete pitch
values from a continuous space. In Western music, $Q$ is a 12-tone equal
temperament (12-TET) scale. In Carnatic music, $Q$ is a 22-\\'{s}ruti system.
In Balinese gamelan, $Q$ is a pelog or slendro tuning. The snap operator is
universal; only the snap points differ.

\begin{lstlisting}[style=lean,caption={Lean 4 formalization of SNAP}]
structure Snap where
  P : Type  -- continuous pitch space
  Q : Finset Float  -- snap points (scale degrees)
  d : Float -> Float -> Float  -- distance metric
  snap : Float -> Float  -- snap operator
  snap_min : forall p, d p (snap p) =
    Finset.min' (Finset.image (fun q => d p q) Q) sorry
\end{lstlisting}

\subsection{\FUNNEL{}: Convergence to Stable Goals}
\label{sec:primitive-funnel}

\begin{definition}[\FUNNEL{}]
A \emph{funnel structure} is a tuple $\mathcal{F} = (S, G, \phi, \tau)$ where:
\begin{itemize}[nosep]
  \item $S$ is a set of musical states (pitch, chord, raga state, etc.),
  \item $G \subset S$ is a set of \emph{goal states} (tonic, sam, downbeat),
  \item $\phi: S \to [0, 1]$ is a \emph{funnel potential} with $\phi(g) = 0$
    for all $g \in G$ and $\phi(s) > 0$ for $s \notin G$,
  \item $\tau: S \to S$ is a \emph{transition function} such that
    $\phi(\tau(s)) \leq \phi(s)$ for all $s \in S$,
\end{itemize}
defining a monotonically decreasing trajectory toward a goal state.
\end{definition}

\FUNNEL{} captures the universal phenomenon of musical resolution: the pull of
the tonic in Western tonality, the approach to \emph{sam} in Carnatic music,
the resolution to a finalis in Arabic maqam, the convergence to a rhythmic
cycle's downbeat in West African drumming.

\begin{lstlisting}[style=lean,caption={Lean 4 formalization of FUNNEL}]
structure Funnel (S : Type) where
  G : Set S  -- goal states
  phi : S -> Float  -- funnel potential
  tau : S -> S  -- transition function
  phi_zero : forall g ∈ G, phi g = 0
  phi_positive : forall s, s ∉ G -> phi s > 0
  monotone : forall s, phi (tau s) <= phi s
\end{lstlisting}

\subsection{\CONSENSUS{}: Multi-Agent Agreement}
\label{sec:primitive-consensus}

\begin{definition}[\CONSENSUS{}]
A \emph{consensus structure} is a tuple $\mathcal{C} = (A, V, \kappa)$ where:
\begin{itemize}[nosep]
  \item $A = \{a_1, \ldots, a_n\}$ is a set of \emph{agents} (musicians,
    voices, parts),
  \item $V$ is a set of possible \emph{values} (pitches, chords, rhythmic
    phases),
  \item $\kappa: A \to V$ is an \emph{agreement function} assigning each
    agent a value,
\end{itemize}
such that there exists an equivalence relation $\sim$ on $V$ where $\kappa(a_i)
\sim \kappa(a_j)$ for all $i, j$ signifies consensus.
\end{definition}

\CONSENSUS{} captures the coordination requirements of ensemble performance:
harmonic agreement in Western chorales, rhythmic alignment in gamelan
interlocking, the mutual tuning of a jazz rhythm section, the unison
convergence of Hindustani musicians on \emph{sam}.

\begin{lstlisting}[style=lean,caption={Lean 4 formalization of CONSENSUS}]
structure Consensus (A V : Type) where
  kappa : A -> V  -- assignment function
  equiv : V -> V -> Prop  -- equivalence relation
  equiv_refl : Reflexive equiv
  equiv_symm : Symmetric equiv
  equiv_trans : Transitive equiv
  agreed : Set A  -- agents that have reached consensus
\end{lstlisting}

\subsection{\LAMAN{}: Structural Rigidity}
\label{sec:primitive-laman}

\begin{definition}[\LAMAN{}]
A \emph{Laman structure} is a pair $\mathcal{L} = (G, \lambda)$ where:
\begin{itemize}[nosep]
  \item $G = (V, E)$ is a graph with $|V| = n$ musical elements and
    $|E| = m$ constraint edges,
  \item $\lambda: E \to \R$ assigns a constraint weight to each edge,
\end{itemize}
such that $G$ satisfies the \emph{Laman condition}: $|E| = 2n - 3$ and every
subset of $k$ vertices spans at least $2k - 3$ edges. This ensures
\emph{minimal rigidity}: the structure has exactly enough constraints to
prevent deformation while remaining flexible enough for musical expression.
\end{definition}

The name \LAMAN{} derives from the Laman graph theorem from rigidity theory
\citep{laman1970graphs}, which we apply to musical structure. In Western
music, the four-voice chorale is a Laman structure: it has exactly enough
contrapuntal constraints to maintain coherence without being overdetermined. In
gamelan, the interlocking kotekan pattern is a Laman structure: it has exactly
enough temporal constraints to maintain rhythmic integrity while allowing
individual players freedom.

\begin{lstlisting}[style=lean,caption={Lean 4 formalization of LAMAN}]
structure Laman where
  V : Finset Nat  -- musical elements (vertices)
  E : Finset (Nat × Nat)  -- constraint edges
  lambda : (Nat × Nat) -> Float  -- edge weights
  edge_count : E.card = 2 * V.card - 3
  laman_cond : forall s ⊆ V,
    (Finset.filter (fun e => e.1 ∈ s ∧ e.2 ∈ s) E).card >= 2 * s.card - 3
\end{lstlisting}

\subsection{\TEMPO{}: Temporal Organization}
\label{sec:primitive-tempo}

\begin{definition}[\TEMPO{}]
A \emph{tempo structure} is a tuple $\mathcal{T} = (\Omega, B, \rho, \gamma)$
where:
\begin{itemize}[nosep]
  \item $\Omega = [0, T)$ is a temporal domain (measured in beats or seconds),
  \item $B = \{b_1, \ldots, b_m\}$ is a set of \emph{beat positions} with
    $b_i \in \Omega$,
  \item $\rho: \Omega \to \R_{> 0}$ is a \emph{tempo function} giving the
    instantaneous rate,
  \item $\gamma: B \to [0, 1]^k$ is a \emph{metric profile} assigning each
    beat a position in a $k$-dimensional metric space,
\end{itemize}
encoding both periodic and non-periodic temporal organization.
\end{definition}

\TEMPO{} captures the universal requirement that music be organized in time.
In Western music, $\rho$ is approximately constant and $\gamma$ encodes a
simple duple or triple meter. In Hindustani music, $\rho$ varies continuously
(rubato) and $\gamma$ encodes the \emph{tala} cycle. In West African drumming,
$\gamma$ encodes multiple simultaneous rhythmic cycles at different
periodicities (polyrhythm).

\begin{lstlisting}[style=lean,caption={Lean 4 formalization of TEMPO}]
structure Tempo where
  omega : Float  -- temporal domain upper bound
  beats : Array Float  -- beat positions
  rho : Float -> Float  -- tempo function
  gamma : Float -> Array Float  -- metric profiles
  rho_pos : forall t, rho t > 0
\end{lstlisting}

% =============================================================================
\section{Cross-Cultural Table: $5 \times 7$ Universality}
\label{sec:table}
% =============================================================================

We now present the central empirical claim of this paper: every one of the five
primitives is instantiated in every one of the seven traditions. Table~\ref
{tab:universal} presents the $5 \times 7 = 35$ cells, each containing a
concrete musical instantiation.

\begin{sidewaystable}[p]
\centering
\caption{Cross-cultural instantiation of the five primitives across seven
traditions. Each cell gives a concrete musical example.}
\label{tab:universal}
\small
\begin{tabular}{@{}l p{3.1cm} p{3.1cm} p{3.1cm} p{3.1cm} p{3.1cm} p{3.1cm} p{3.1cm}@{}}
\toprule
\textbf{Primitive} & \textbf{Western Tonal} & \textbf{Carnatic} & \textbf{Hindustani} & \textbf{Jazz} & \textbf{W.\ African} & \textbf{Arabic Maqam} & \textbf{Balinese Gamelan} \\
\midrule
\textbf{\SNAP{}} &
12-TET chromatic scale; pitch classes $\{0, 1, \ldots, 11\}$ &
22-\\'{s}ruti system; microtonal divisions of the octave &
12 swara with microtonal shading (shruti) &
Chromatic + blues notes; pitch-bend as soft snap &
7-tone heptatonic scales with flexible intonation &
24-tone equal division; quarter-tone snap points &
Pelog (7-tone) and slendro (5-tone) non-equidistant scales \\
\midrule
\textbf{\FUNNEL{}} &
Dominant--tonic resolution; $V \to I$ cadence &
Approach to sam; vadi--samvadi convergence &
Convergence to shadja (sa); nyas tones &
ii--V--I progression; tritone resolution &
Call-and-response convergence to rhythmic cycle &
Finalis resolution; sayr (melodic path) to tonic &
Gong cycle convergence; approach to strongest beat \\
\midrule
\textbf{\CONSENSUS{}} &
SATB chorale voicing; harmonic agreement &
Unison at sam; multiple performers converge &
Drone agreement (tanpura + soloist) &
Rhythm section lock; walking bass + ride cymbal &
Interlocking patterns (bell + drum agreement) &
Takht ensemble tuning; agreement on maqam &
Kotekan interlock; two-part agreement \\
\midrule
\textbf{\LAMAN{}} &
4-voice species counterpoint; $2n - 3$ constraints &
Varna structure; phrase-level rigidity &
Khayal form constraints; elastic Laman &
Chord-scale rigidity; ii--V--I as rigid frame &
Bell pattern as Laman graph; temporal rigidity &
Waslah structure; suite-level minimal rigidity &
Colotomic structure; gong cycle as rigid frame \\
\midrule
\textbf{\TEMPO{}} &
$4/4$ or $3/4$ meter; isochronous beats &
Tala cycles (Adi: $4{+}2{+}2$); asymmetric &
Tintal (16 beats), Jhaptal (10 beats) &
Swing feel; asymmetric beat subdivision &
Polyrhythmic layers; 12/8 bell patterns &
Iqa'at cycles; maqsum (4/4), samai (10/8) &
Gong cycles; tempo doubles at structure points \\
\bottomrule
\end{tabular}
\end{sidewaystable}

\subsection{Reading the Table}

Each cell in Table~\ref{tab:universal} represents a specific instance of a
primitive within a tradition. The universality claim is not that the
instantiations are identical---they manifestly are not---but that they are
instances of the \emph{same mathematical structure}. The Western dominant--tonic
cadence and the Carnatic approach to \emph{sam} are both instances of
\FUNNEL{}: they are monotonically decreasing trajectories toward a goal state,
differing only in the choice of state space $S$, goal set $G$, and potential
function $\phi$.

\subsection{Scale and Tuning Taxonomy}

Our framework encompasses 36 distinct scales across 7 tuning systems:

\begin{itemize}[nosep]
  \item \textbf{Western:} major, natural minor, harmonic minor, melodic
    minor, Dorian, Phrygian, Lydian, Mixolydian, Locrian, whole-tone,
    chromatic, blues (12 scales, 12-TET).
  \item \textbf{Carnatic:} the 72 \emph{melakarta} ragas are generated by a
    22-\\'{s}ruti snap; we enumerate 8 representative scales.
  \item \textbf{Hindustani:} 10 principal \emph{thaat} with microtonal
    variants.
  \item \textbf{Jazz:} bebop scales, altered scales, pentatonic, octatonic
    (6 scales, 12-TET with blues microtones).
  \item \textbf{West African:} heptatonic and pentatonic scales with flexible
    intonation.
  \item \textbf{Arabic:} maqam scales in 24-TEDO (Rast, Bayati, Hijaz,
    Kurd, Saba, Ajam, Nakriz, Sikah; 8 representative scales).
  \item \textbf{Balinese:} pelog and slendro tunings with regional variants
    (4 scales).
\end{itemize}

Similarly, 26 rhythmic patterns span the traditions: Western meters ($2/4$,
$3/4$, $4/4$, $6/8$, $5/4$, $7/8$), Carnatic talas (Adi, Rupaka,
Misra Chapu, Khanda Chapu, Dhruva, Matya), Hindustani talas (Tintal, Jhaptal,
Ektal, Rupak, Keherwa), jazz comping patterns, West African bell patterns
(12/8 standard, Ewe agbadza, Yoruba sakara), Arabic iqa'at (Maqsum, Baladi,
Ayd, Samai Thaqil, Masmudi Kabir, Dawr Hindy), and Balinese gong cycles.

% =============================================================================
\section{Implementation}
\label{sec:implementation}
% =============================================================================

We provide reference implementations of Constraint Theory in seven programming
languages, each chosen for a specific purpose:

\begin{enumerate}[nosep]
  \item \textbf{Haskell:} The primary implementation, leveraging algebraic data
    types and type classes for the mathematical structures.
  \item \textbf{Python:} A practical implementation for data science workflows,
    with NumPy arrays for pitch and rhythm manipulation.
  \item \textbf{Rust:} A high-performance implementation for real-time
    generation, targeting embedded systems and live coding.
  \item \textbf{TypeScript:} A web-based implementation for browser-based
    composition tools and interactive visualizations.
  \item \textbf{Lean 4:} The formal verification implementation (\S\ref
    {sec:verification}).
  \item \textbf{Julia:} A scientific computing implementation for large-scale
    analysis and optimization.
  \item \textbf{Pure Data:} A visual programming implementation for real-time
    audio processing and live performance.
\end{enumerate}

\subsection{Haskell Core}

The Haskell implementation uses type classes to encode the five primitives:

\begin{lstlisting}[style=haskell,caption={Haskell type classes for the five primitives}]
class Snap s where
  type PitchSpace s :: Type
  snapPoints :: s -> [Float]
  snap :: s -> Float -> Float

class Funnel f where
  type State f :: Type
  goals :: f -> Set (State f)
  potential :: f -> State f -> Float
  transition :: f -> State f -> State f)

class Consensus c where
  type Agent c :: Type
  type Value c :: Type
  agree :: c -> Map (Agent c) (Value c)
  isConsensus :: c -> Bool

class Laman l where
  vertices :: l -> Int
  edges :: l -> [(Int, Int, Float)]
  isRigid :: l -> Bool

class Tempo t where
  beats :: t -> [Float]
  tempoFn :: t -> Float -> Float
  metricProfile :: t -> Float -> [Float]
\end{lstlisting}

\subsection{Cross-Traditional Instantiation}

Each tradition is modeled as a record bundling the five primitives:

\begin{lstlisting}[style=haskell,caption={Tradition as a bundle of primitives}]
data Tradition = Tradition
  { snapInst     :: SomeSnap
  , funnelInst   :: SomeFunnel
  , consensusInst :: SomeConsensus
  , lamanInst    :: SomeLaman
  , tempoInst    :: SomeTempo
  , traditionName :: String
  }

westernTonal, carnatic, hindustani :: Tradition
jazz, westAfrican, arabicMaqam, balinese :: Tradition
\end{lstlisting}

% =============================================================================
\section{\SNAP{} Universality}
\label{sec:snap}
% =============================================================================

\begin{proposition}[Universality of \SNAP{}]
Every musical tradition with a defined pitch vocabulary has a snap structure.
\end{proposition}

\begin{proof}[Proof sketch]
By definition, any tradition that uses a finite set of pitch classes defines
snap points $Q$. The pitch space $P$ is $\R^+$ (frequency) or $[0, 1200)$
(cents modulo octave). The distance $d$ is the standard metric on $\R$. Since
$Q$ is finite and $P$ is compact modulo octave equivalence, the snap operator
$\sigma$ exists uniquely for all $p$ not equidistant between two snap points.
At such boundary points, a tie-breaking rule (upward or downward snap) resolves
uniqueness. Since all seven traditions in our study use finite pitch sets,
\SNAP{} is universal among them.
\end{proof}

The key insight is that snap points differ across traditions, but the
\emph{structure} of quantization is the same. The 22-\\'{s}ruti system of
Carnatic music and the 12-TET system of Western music are both instances of
$\sigma: \R \to Q$ for different choices of $Q$. Even the continuous pitch
inflections of Hindustani music (meend) can be modeled as sequences of snap
operations over a time-varying $Q(t)$.

\subsection{Multi-Scale Snap}

A single piece may involve multiple snap structures simultaneously. In jazz, a
performer may snap to the chord tones of the current harmony (a time-varying
$Q(t)$), the underlying scale (a static $Q$), and the blues scale (another
static $Q$), navigating between them fluidly. This is formalized as a
\emph{multi-snap structure}:

\begin{definition}[Multi-Snap]
A \emph{multi-snap structure} is a family $\{\mathcal{S}_i\}_{i \in I}$ of
snap structures on the same pitch space $P$, together with a weight function
$w: I \times P \to \R_{\geq 0}$, such that the composite snap operator is:
\[
  \sigma^*(p) = \arg\min_{q \in \bigcup_i Q_i} \sum_i w(i, p) \cdot d(p, q).
\]
\end{definition}

Multi-snap structures arise naturally in many traditions. In Arabic maqam,
the performer navigates between the scale degrees of the current maqam and
those of closely related maqamat during modulation. The weight function $w$
shifts from the current maqam's snap points to the target maqam's snap points,
creating a smooth transition that is perceived as a modulation rather than a
sudden change. In West African vocal music, the singer may snap to different
scale degrees depending on the linguistic tone of the text being sung---a
remarkable example of multi-snap where the weight function is determined by
language rather than harmony.

The multi-snap framework also captures the phenomenon of ``blue notes'' in
jazz and blues. A blue note is simultaneously a snap to the diatonic scale
(with low weight) and a snap to the blues scale (with high weight). The
resulting composite snap produces a pitch that is not a member of either
scale but is attracted to both---a hallmark of the African-American musical
tradition's creative negotiation between multiple pitch systems.

% =============================================================================
\section{\FUNNEL{} Universality}
\label{sec:funnel-sec}
% =============================================================================

\begin{proposition}[Universality of \FUNNEL{}]
Every musical tradition with defined melodic goals has a funnel structure.
\end{proposition}

\begin{proof}[Proof sketch]
Define the state space $S$ as the set of possible pitch-durations or
harmonic-rhythmic configurations. Define $G$ as the set of ``resting''
states: the tonic, the finalis, the sam, the downbeat. The potential $\phi$
measures distance from the current state to the nearest goal state; by
construction, $\phi(g) = 0$ for $g \in G$. The transition function $\tau$
encodes the next-state relation of the tradition's grammar (voice leading
rules, raga grammar, iqa' fill patterns). Since musical practice universally
involves departure from and return to resting states, $\phi(\tau(s)) \leq
\phi(s)$ along ``normative'' paths. The funnel structure thus exists.
\end{proof}

The \FUNNEL{} primitive is closely related to the concept of
\emph{expectation} in music cognition \citep{meyer1956emotion,huron2006sweet}.
The funnel potential $\phi$ encodes the listener's expectation of resolution;
the transition function $\tau$ encodes the composer's or performer's
choices for satisfying (or deliberately frustrating) that expectation.

\subsection{Funnel Topology}

The topology of the funnel---the structure of goal states and the paths
leading to them---varies across traditions and provides a powerful
characterization of each tradition's melodic style.

In Western tonal music, the funnel topology is \emph{hierarchical}: there are
primary goals (the tonic triad), secondary goals (the dominant, predominant
harmonies), and tertiary goals (passing tones, neighboring tones). The
hierarchy corresponds to the hierarchy of harmonic function (tonic, dominant,
predominant). The funnel potential $\phi$ reflects this hierarchy: a phrase
ending on the dominant has a higher $\phi$ than one ending on the tonic, and a
phrase in the middle of a sequence has a higher $\phi$ than one at a
structural cadence.

In Carnatic music, the funnel topology is \emph{cyclical}: the goal is
\emph{sam} (beat 1 of the tala cycle), and the funnel wraps around the
rhythmic cycle. The potential $\phi$ decreases as the performer approaches
sam from any direction, creating a cyclical pattern of tension and release
that repeats with each avartana (cycle). This cyclical topology is shared with
West African drumming, where the bell pattern creates a cyclical funnel toward
the main beat.

In Hindustani music, the funnel topology is \emph{layered}: there are
simultaneous funnels operating at different timescales. At the micro level,
each phrase has its own funnel (approaching a nyas tone). At the meso level,
the section has a funnel (approaching the climax). At the macro level, the
entire performance has a funnel (approaching the final resolution). These
nested funnels interact, creating a rich tension landscape that unfolds over
the course of a performance.

In jazz, the funnel topology is \emph{progressive}: the ii--V--I progression
creates a sequence of funnels, each resolving into the next. The overall
trajectory is progressive (moving through the harmonic form) rather than
cyclical (returning to the same point). However, the jazz \emph{blues} form
adds a cyclical component: the 12-bar form creates a macro-level funnel that
returns to the tonic at bar 1, while the internal harmonies create micro-level
progressive funnels.

\subsection{Funnel Disruption}

One of the most expressive techniques across traditions is the deliberate
disruption of the funnel: creating an expectation of resolution and then
denying it. In Western music, this is the deceptive cadence ($V \to vi$
instead of $V \to I$). In Hindustani music, it is the deliberate avoidance
of the shadja (tonic) at an expected resolution point. In jazz, it is the
tritone substitution ($\mathrm{ii} \to \mathrm{\flat II} \to I$ instead of
$\mathrm{ii} \to V \to I$).

Formally, funnel disruption occurs when the transition function $\tau$
produces a state with $\phi(\tau(s)) > \phi(s)$---an \emph{increase} in
potential rather than the expected decrease. The magnitude of the disruption
is $\Delta\phi = \phi(\tau(s)) - \phi(s) > 0$. This quantity measures the
``surprise'' of the disruption. The aesthetic effect depends on the subsequent
resolution: a large disruption followed by a strong resolution creates a
powerful emotional arc (the ``deceptive cadence effect'').

We hypothesize that the distribution of $\Delta\phi$ values in a performance
correlates with perceived emotional intensity and that the distribution differs
across traditions in predictable ways. Jazz, for instance, is expected to have
a wider distribution of $\Delta\phi$ than Western classical, reflecting the
tradition's greater tolerance for harmonic surprise.

% =============================================================================
\section{\CONSENSUS{} Universality}
\label{sec:consensus-sec}
% =============================================================================

\begin{proposition}[Universality of \CONSENSUS{}]
Every ensemble musical tradition has a consensus structure.
\end{proposition}

The proof follows from the observation that any ensemble performance requires
agreement on at least one musical parameter: pitch (tuning), rhythm (tempo
synchronization), or form (structural boundaries). Even freely improvised
music requires temporal consensus (beginning and ending together).

\CONSENSUS{} is formalized as a distributed consensus problem, connecting
music theory to distributed computing \citep{lamport1982consensus}. The
analogy is deep: musicians in an ensemble are agents that must reach agreement
despite latency (the speed of sound), partial information (you cannot hear
every player equally well), and Byzantine failures (a player might make a
mistake).

In Balinese gamelan, the \emph{kotekan} (interlocking figuration) requires two
players to jointly produce a single composite melody, with each player
contributing alternate notes. This is a consensus problem where the agreement
relation is temporal: the two parts must interleave correctly to form the
composite melody. The remarkable precision of gamelan coordination
\citep{tenzer2000gamelan} suggests that musicians have developed embodied
solutions to consensus that rival engineered systems.

In Arabic music, the \emph{takht} ensemble must agree on the current maqam
and modulations within it. The agreement is dynamic: as the performance
progresses, the ensemble collectively navigates the maqam's sayr (melodic
path), and any deviation by one player must be matched by the others to
maintain coherence. This is a distributed consensus problem with a time-varying
goal: the target maqam changes as the piece modulates.

Jazz presents perhaps the richest consensus problem in our framework. A jazz
ensemble must agree on multiple parameters simultaneously: harmony (which
chord-scale to use), rhythm (subdivision and swing feel), form (location in
the tune's structure), and dynamics. Crucially, jazz consensus is
\emph{emergent}: there is no conductor, and the agreement arises from mutual
listening and response. The tempo may shift slightly (elastic tempo), the
harmony may be reinterpreted (reharmonization), and the form may be altered
(tagging the ending, extending solos), yet the ensemble maintains consensus
through continuous micro-adjustments. This is analogous to the \emph{flocking}
behavior studied in multi-agent systems \citep{olfati2004consensus}, where
global coherence emerges from local interactions.

\subsection{Consensus Dynamics}

We model consensus as a dynamical system. Let $x_i(t) \in V$ be the value of
agent $a_i$ at time $t$. The consensus dynamics are:

\begin{equation}
  \dot{x}_i(t) = \sum_{j \in N(i)} w_{ij} \bigl(x_j(t) - x_i(t)\bigr),
  \label{eq:consensus-dynamics}
\end{equation}

where $N(i)$ is the set of agents audible to $a_i$ and $w_{ij}$ is an
attention weight. This is the standard linear consensus protocol from
multi-agent systems theory \citep{olfati2004consensus}. Convergence to
consensus is guaranteed when the communication graph (who can hear whom) is
connected, which is always the case in a typical ensemble setup.

% =============================================================================
\section{\LAMAN{} Universality}
\label{sec:laman-sec}
% =============================================================================

\begin{proposition}[Universality of \LAMAN{}]
Every musical tradition with identifiable structural constraints has a Laman
structure (possibly elastic).
\end{proposition}

The connection between musical structure and rigidity theory is, to our
knowledge, new. The key observation is that musical constraints (voice-leading
rules, counterpoint rules, raga grammar) are \emph{distance constraints} on a
graph: they specify that certain pairs of musical elements must maintain a
certain relationship (interval, temporal offset, metric position). A
minimally rigid structure has exactly enough constraints to determine the
structure without overdetermining it.

Rigidity theory, originating in the study of mechanical structures
\citep{laman1970graphs}, provides precisely the right language for this
phenomenon. A Laman graph on $n$ vertices has $2n - 3$ edges and every
subgraph on $k$ vertices spans at least $2k - 3$ edges. This ensures that the
structure is rigid (it cannot deform freely) but not overconstrained (no edge
is redundant). In music, this corresponds to the ideal balance between
structure (enough rules to ensure coherence) and freedom (enough flexibility
for expression).

The analogy extends further. In rigidity theory, a structure that has more
than $2n - 3$ edges is \emph{overconstrained}: some edges are redundant, and
removing them does not affect rigidity. In music, an overconstrained structure
corresponds to a style that has more rules than necessary---a pedantic
realization of species counterpoint, perhaps, where every interval is
prescribed and no artistic choice remains. Conversely, a structure with fewer
than $2n - 3$ edges is \emph{floppy}: it can deform freely. In music, this
corresponds to free improvisation with insufficient framework to maintain
coherence.

In Western four-voice chorale harmony, the constraint graph has:
\begin{itemize}[nosep]
  \item Vertices: the four voices (soprano, alto, tenor, bass).
  \item Edges: interval constraints between each pair of voices at each time
    step, plus melodic constraints (smooth voice leading) within each voice
    across time steps.
\end{itemize}

The Laman condition $|E| = 2|V| - 3$ ensures that the harmony is fully
determined by the constraints (no extra degrees of freedom) but not
overconstrained (some flexibility remains for artistic choice).

\subsection{Elastic Laman}

Real musical structures are not perfectly rigid; they allow controlled
deformation. We introduce \emph{elastic Laman} structures, where the edge
weights $\lambda(e)$ are treated as spring constants rather than hard
constraints:

\begin{equation}
  E(\mathbf{x}) = \sum_{e \in E} \lambda(e) \bigl(d(\mathbf{x}_i, \mathbf{x}_j)
  - \ell_e\bigr)^2,
  \label{eq:elastic-laman}
\end{equation}

where $\ell_e$ is the rest length (desired interval or offset) and $d$ is the
distance function. The musical structure is the configuration $\mathbf{x}^*$
that minimizes $E$. This captures the flexibility of real music: rules are
guidelines, not absolute constraints, and the ``best'' musical structure
balances competing demands.

In Hindustani music, the \emph{raga} grammar specifies allowed and forbidden
melodic movements. An elastic Laman structure captures this by assigning high
$\lambda$ values to forbidden movements (strong penalty) and lower values to
preferred movements (weak attraction). The performer navigates the energy
landscape $E$, choosing paths that honor the constraints while expressing
creativity.

\subsection{Laman Structures Across Traditions}

We briefly sketch the Laman structure for each tradition:

\begin{itemize}[nosep]
  \item \textbf{Western tonal:} The four-voice SATB chorale is the canonical
    example. The constraint graph connects each pair of voices with interval
    constraints (parallel fifths forbidden, dissonances prepared and resolved)
    and each voice with its own temporal predecessor (smooth voice leading).
    The Laman condition ensures that exactly one note in each voice is free
    at each time step (the others are determined by constraints), giving the
    composer exactly the right amount of freedom.
  \item \textbf{Carnatic:} The varnam structure provides a Laman graph at
    the phrase level. Each phrase (prayogam) has constraints connecting it to
    its predecessor (sequential constraint), the tala cycle (temporal
    constraint), and the raga's vadi--samvadi (pitch constraint). The Laman
    condition ensures that each phrase is determined by the raga, the tala,
    and the preceding phrase, leaving room for the performer's
    manodharma (imagination).
  \item \textbf{Jazz:} The chord-scale relationship provides a Laman graph
    at the harmonic level. Each chord is connected to its scale (pitch
    constraint), its predecessor and successor in the progression (voice
    leading constraint), and the metrical position (temporal constraint).
    The Laman condition ensures that the improviser has enough constraints to
    maintain coherence but enough freedom for spontaneous creation.
  \item \textbf{West African:} The bell pattern provides a temporal Laman
    graph. Each stroke is connected to its neighbors (sequential constraint)
    and to the underlying pulse (metric constraint). The Laman condition
    ensures that the pattern is recognizably the same across repetitions
    but allows micro-timing variation.
  \item \textbf{Arabic maqam:} The waslah (suite) structure provides a
    Laman graph at the formal level. Each section is connected to its
    neighbors (sequential constraint), the underlying maqam (pitch constraint),
    and the iqa' (rhythmic constraint). The Laman condition ensures formal
    coherence across a long, multi-section performance.
  \item \textbf{Balinese gamelan:} The colotomic structure (gong cycle)
    provides a temporal Laman graph. Each gong stroke is connected to the
    preceding and following strokes (sequential constraint), the underlying
    tempo (metric constraint), and the dance or dramatic context (functional
    constraint). The Laman condition ensures that the ensemble maintains
    temporal coherence across the complex, layered texture.
\end{itemize}

% =============================================================================
\section{\TEMPO{} Universality}
\label{sec:tempo}
% =============================================================================

\begin{proposition}[Universality of \TEMPO{}]
Every musical tradition with perceivable temporal organization has a tempo
structure.
\end{proposition}

The proof is trivial: any music that unfolds in time can be described by a set
of event times, which define $B$. The tempo function $\rho$ may be constant
(Western art music at steady tempo), periodic (rubato in Romantic piano),
or more complex (the accelerating tempos of gamelan). The metric profile
$\gamma$ captures the hierarchical structure of meter.

Our framework encompasses 26 rhythmic patterns, organized by the type of
\TEMPO{} structure they employ:

\begin{itemize}[nosep]
  \item \textbf{Isochronous:} Constant $\rho$, simple $\gamma$ (Western
    $4/4$, $3/4$; jazz walking bass).
  \item \textbf{Hierarchical:} Constant $\rho$, nested $\gamma$ (Carnatic
    Adi tala with matra--akshara--avartana hierarchy).
  \item \textbf{Polyrhythmic:} Constant $\rho$, multiple $\gamma$ with
    coprime periods (West African 3:2, 4:3, 7:4).
  \item \textbf{Additive:} Constant $\rho$, non-isochronous $\gamma$
    (Brahmsian $5/4$ as $2+3$, Balkan $7/8$ as $2+2+3$).
  \item \textbf{Accelerating:} Increasing $\rho$ (gamelan tempo doubling;
    Hindustani ati-lay to drut-lay transitions).
  \item \textbf{Free:} Undefined $\rho$, event-driven $B$ (alap, unmetered
    taqsim; modeled as $\TEMPO{}$ with $\rho \to 0$).
\end{itemize}

\subsection{The 26 Rhythmic Patterns}

For completeness, we enumerate the 26 rhythmic patterns covered by our
framework:

\noindent\textbf{Western (6):} $2/4$, $3/4$, $4/4$, $6/8$, $5/4$, $7/8$.

\noindent\textbf{Carnatic (6):} Adi ($4{+}2{+}2 = 8$ beats), Rupaka ($2{+}4 = 6$
beats, counting from the sam), Misra Chapu ($3{+}2{+}2 = 7$ beats), Khanda
Chapu ($2{+}1{+}2 = 5$ beats), Dhruva ($1{+}1{+}2{+}2 = 6$ matras per
cycle), Matya ($1{+}2{+}1 = 4$ matras per cycle).

\noindent\textbf{Hindustani (5):} Tintal (16 beats, $4{+}4{+}4{+}4$), Jhaptal
(10 beats, $2{+}3{+}2{+}3$), Ektal (12 beats, $2{+}2{+}2{+}2{+}2{+}2$),
Rupak (7 beats, $3{+}2{+}2$), Keherwa (8 beats, $4{+}4$).

\noindent\textbf{Jazz (2):} Swing 4/4 (with triplet subdivision), Jazz waltz
(3/4 with swing).

\noindent\textbf{West African (3):} 12/8 standard bell pattern (Ewe), Agbadza
bell pattern, Sakara bell pattern (Yoruba).

\noindent\textbf{Arabic (3):} Maqsum ($2{+}2{+}2{+}2$ in 4/4), Samai Thaqil
($3{+}2{+}2{+}3$ in 10/8), Masmudi Kabir ($4{+}4{+}4{+}4$ with specific
accent patterns).

\noindent\textbf{Balinese (1):} Gongan cycle (variable length, typically
$4$, $8$, $16$, or $32$ beats with colotomic punctuation).

These 26 patterns represent the major rhythmic organizing principles found
across the seven traditions. Each can be fully described as a \TEMPO{}
structure $\mathcal{T} = (\Omega, B, \rho, \gamma)$, demonstrating the
universality of the primitive.

% =============================================================================
\section{Soft Snap $\eps$ and the Goldilocks Hypothesis}
\label{sec:softsnap}
% =============================================================================

The snap operator $\sigma$ is binary: a pitch $p$ snaps to the nearest
$q \in Q$ or it does not. Real music is not so clean. Jazz musicians play
between the cracks of the piano; Hindustani singers slide through
microtonal inflections; blues guitarists bend notes to express emotion. We
model this with a \emph{soft snap} operator parameterized by $\eps \geq 0$.

\begin{definition}[Soft Snap]
The \emph{soft snap} operator $\sigma_\eps: \R \to \R$ with parameter
$\eps \geq 0$ is:
\[
  \sigma_\eps(p) = \frac{\displaystyle\sum_{q \in Q} q \cdot \exp\!\left(
  -\frac{d(p, q)^2}{2\eps^2}\right)}{\displaystyle\sum_{q \in Q} \exp\!\left(
  -\frac{d(p, q)^2}{2\eps^2}\right)},
\]
where $\sigma_0 = \sigma$ (hard snap) and $\sigma_\infty = \mathrm{id}$
(identity, no snap).
\end{definition}

The parameter $\eps$ controls the ``sharpness'' of quantization:
\begin{itemize}[nosep]
  \item $\eps \to 0$: Hard snap. The pitch is quantized to the nearest scale
    degree. Corresponds to strict adherence to tradition.
  \item $\eps$ small: Near-snap. The pitch is close to a scale degree but
    with microtonal shading. Corresponds to expressive intonation within a
    tradition (Carnatic gamakas, blues bends).
  \item $\eps$ moderate: Weak snap. The pitch is attracted to scale degrees
    but explores the space between them. Corresponds to stylistic innovation
    (chromatic jazz lines, raga fusion).
  \item $\eps \to \infty$: No snap. The pitch is free. Corresponds to
    avant-garde or free improvisation outside any tradition.
\end{itemize}

\subsection{The Goldilocks Hypothesis}

\begin{conjecture}[Goldilocks Hypothesis]
For any tradition $T$ and any musical piece $M$ generated within $T$'s
constraint framework, there exists an interval $(\eps_{\min}, \eps_{\max})$
such that $M$ is perceived as aesthetically compelling by informed listeners if
and only if $\eps_M \in (\eps_{\min}, \eps_{\max})$.
\end{conjecture}

The Goldilocks Hypothesis states that the best music is neither too
constrained ($\eps$ too small: rigid, predictable, boring) nor too free
($\eps$ too large: random, incoherent, unmusical), but ``just right.'' This
echoes Berlyne's inverted-U theory of aesthetic preference
\citep{berlyne1971aesthetics} and recent work on ``sweet spots'' in creative
AI \citep{boden2004creative}.

We hypothesize that the Goldilocks interval differs across traditions:
\begin{itemize}[nosep]
  \item \textbf{Western tonal:} Narrow interval; Western art music tolerates
    less deviation from the tempered scale and prescribed voice leading.
    The 12-TET snap is hard, and pitch deviation is heard as out-of-tune.
  \item \textbf{Jazz:} Wider interval; jazz values individual expression and
    ``playing outside.'' Blue notes are \emph{expected} deviations from the
    nominal scale, and the tradition's aesthetic rewards controlled
    transgression.
  \item \textbf{Hindustani:} Moderate interval; the raga framework is strict
    in its prescribed notes and forbidden movements, but gamakas (ornamental
    oscillations) allow wide microtonal variation \emph{within} the
    constraints. The tradition demands both fidelity and fluidity.
  \item \textbf{Carnatic:} Moderate-narrow interval; while the 22-\'{s}ruti
    system provides a finer snap grid, the grammar of gamakas and the
    requirements of kalpanaswara improvisation demand precision within the
    microtonal framework.
  \item \textbf{West African:} Wide interval in pitch, narrow in rhythm;
    rhythmic precision is paramount (the bell pattern must be exact), but
    pitch is flexible and vocal intonation is an expressive parameter.
  \item \textbf{Arabic maqam:} Moderate interval; the 24-TEDO system provides
    a finer grid than 12-TET, and the tradition's microtonal ornaments
    (quarter tones, inflections) are mandatory rather than optional. The
    Goldilocks band encompasses both the precision of the maqam's tonic and
    the expressive freedom of sayr.
  \item \textbf{Balinese gamelan:} Narrow interval; the ensemble must align
    precisely on the fixed-pitch metallophones. The tuning itself is
    tradition-specific (each gamelan has its own tuning), but within a given
    ensemble, the snap is extremely hard.
\end{itemize}

These predictions can be tested empirically. For each tradition, we can
construct a corpus of ``canonical'' performances, measure the distribution of
$\eps$ values across the corpus (by computing the deviation from the nearest
snap point for each pitch), and compare the resulting distribution to the
Goldilocks prediction. We expect to find that the distribution is unimodal
with a peak in the predicted interval.

The Goldilocks Hypothesis has a natural interpretation in terms of information
theory. The soft snap $\sigma_\eps$ defines a channel from the continuous
pitch space to the quantized output. As $\eps$ varies, the mutual information
between input and output changes. The Goldilocks interval corresponds to a
range of mutual information that is neither too low (too much information
lost) nor too high (too little information lost), consistent with
\emph{optimal information transport} theories of aesthetic perception
\citep{birkhoff1933aesthetic}.

\subsection{Soft Funnel, Soft Consensus, Soft Laman}

The soft snap paradigm extends naturally to the other primitives:

\begin{itemize}[nosep]
  \item \textbf{Soft Funnel:} $\phi_\eps(s) = (1 - \eps) \cdot \phi(s) + \eps
    \cdot \phi_{\mathrm{free}}(s)$, interpolating between the tradition's
    funnel potential and a free (flat) potential.
  \item \textbf{Soft Consensus:} The agreement relation $\sim_\eps$ relaxes
    to allow ``approximate'' agreement within tolerance $\eps$.
  \item \textbf{Soft Laman:} Edge weights become $\lambda_\eps(e) = (1-\eps)
    \cdot \lambda(e) + \eps \cdot 0$, interpolating toward a disconnected
    (unconstrained) graph.
  \item \textbf{Soft Tempo:} The beat grid becomes probabilistic with
    tolerance $\eps$ around each nominal beat position.
\end{itemize}

A unified soft constraint framework is parameterized by a vector
$\boldsymbol{\eps} = (\eps_s, \eps_f, \eps_c, \eps_l, \eps_t) \in \R_{\geq
0}^5$, controlling the ``traditionness'' of each primitive independently.
This vector provides a five-dimensional space for exploring the relationship
between tradition and innovation.

% =============================================================================
\section{Cross-Cultural Composition via Hyperbolic Geometry}
\label{sec:hyperbolic}
% =============================================================================

The five primitives provide a common language for describing different musical
traditions. But how do we \emph{blend} traditions? A naive approach---mixing
snap points, for instance---produces chromatic mash-ups with no coherence. We
need a principled method for interpolation.

We propose that musical traditions occupy distinct regions of a \emph{constraint
manifold} $\mathcal{M}$, and that cross-cultural blending corresponds to
geodesic interpolation in this manifold. The natural geometry for this manifold
is \emph{hyperbolic}.

\subsection{The Constraint Manifold}

\begin{definition}[Constraint Manifold]
The \emph{constraint manifold} $\mathcal{M}$ is the product space:
\[
  \mathcal{M} = \mathcal{S} \times \mathcal{F} \times \mathcal{C} \times
  \mathcal{L} \times \mathcal{T},
\]
where each factor is the space of possible configurations of the corresponding
primitive. A \emph{tradition} is a region $T_i \subset \mathcal{M}$ defined by
the specific parameter values of that tradition.
\end{definition}

Each tradition $T_i$ is not a point in $\mathcal{M}$ but a region, because a
tradition encompasses many possible pieces, performances, and interpretations.
The region has a natural center (the ``prototypical'' configuration) and a
boundary (the limits of what is recognizable as that tradition).

\subsection{Hyperbolic Structure}

We equip $\mathcal{M}$ with a hyperbolic metric $g_H$ modeled on the
Poincar\'{e} ball. The intuition is that the ``center'' of a tradition (where
$\eps = 0$) is at the boundary of hyperbolic space, and the free space
($\eps = \infty$) is at the origin. This geometry captures the empirical
observation that small departures from a tradition are easy (near the
boundary, distances are large) but large departures are ``far apart'' in a
meaningful sense.

Formally, let $\mathbf{p}_i \in \mathcal{M}$ be the center of tradition $T_i$,
represented as a point in the Poincar\'{e} ball $B^n = \{\mathbf{x} \in \R^n :
\|\mathbf{x}\| < 1\}$. The hyperbolic distance between two traditions is:

\begin{equation}
  d_H(\mathbf{p}_i, \mathbf{p}_j) = \mathrm{arcosh}\!\left(1 + 2
  \frac{\|\mathbf{p}_i - \mathbf{p}_j\|^2}{(1 - \|\mathbf{p}_i\|^2)(1 -
  \|\mathbf{p}_j\|^2)}\right).
  \label{eq:hyperbolic-distance}
\end{equation}

\subsection{Geodesic Blending}

Given two traditions $T_i$ and $T_j$, a \emph{blend} is a piece whose
constraint parameters lie along the geodesic $\gamma: [0, 1] \to \mathcal{M}$
connecting $\mathbf{p}_i$ and $\mathbf{p}_j$:

\begin{equation}
  \gamma(t) = \frac{(1-t)\mathbf{p}_i / (1 - \|\mathbf{p}_i\|^2) +
  t \mathbf{p}_j / (1 - \|\mathbf{p}_j\|^2)}
  {1 + \|\cdot\|^2_{H}} \cdot \frac{1}{\text{norm}},
  \label{eq:geodesic}
\end{equation}

where the normalization ensures $\gamma(t)$ remains in $B^n$. The blend at
$t = 0$ is purely $T_i$; at $t = 1$ purely $T_j$; and at intermediate values
it is a smooth hybrid.

\subsection{Examples of Geodesic Blending}

\begin{itemize}[nosep]
  \item \textbf{Jazz--Carnatic blend ($t = 0.5$):} The snap points are a
    weighted combination of 12-TET and 22-\\'{s}ruti, producing a
    microtonally enriched jazz harmony. The funnel combines ii--V--I
    resolution with gamaka-based approach to sam. The consensus structure
    allows both jazz rhythm-section coordination and Carnatic unison
    convergence.
  \item \textbf{West African--Balinese blend ($t = 0.5$):} The tempo
    structure combines polyrhythmic layers with gong-cycle colotomy. The
    Laman structure combines bell-pattern rigidity with kotekan interlocking.
    The consensus structure enables both rhythmic alignment and melodic
    interlock.
  \item \textbf{Arabic--Hindustani blend ($t = 0.5$):} The snap points
    combine 24-TEDO quarter tones with microtonal shruti shading. The funnel
    combines maqam sayr with raga vadi--samvadi convergence.
\end{itemize}

The hyperbolic geometry ensures that ``nearby'' traditions (those sharing many
features) blend more naturally than ``distant'' ones. This reflects the
empirical observation that, e.g., Hindustani and Carnatic music blend more
easily than Hindustani and Balinese music, because the former share more
structural features.

\subsection{Curvature and Cultural Distance}

The curvature of the constraint manifold encodes cultural distance. In a flat
(Euclidean) geometry, the shortest path between two traditions is a straight
line---a simple weighted average of parameters. In a hyperbolic geometry, the
geodesic is curved, and the intermediate points are ``pulled'' toward the
origin (the free space). This means that a blend of two distant traditions
passes through a region of greater freedom before converging on the target.

This has a natural musical interpretation: blending distant traditions
requires passing through a ``free zone'' where neither tradition's constraints
fully apply. This free zone is where the most creative (and potentially the
most interesting) music emerges, but also where the risk of incoherence is
greatest. The Goldilocks Hypothesis applies here in a new guise: the best
blends are those that navigate the free zone without losing coherence.

The hyperbolic structure also explains why certain tradition pairs are
``natural allies'' in fusion music. Jazz and Indian classical music, for
example, have a long history of successful fusion (John McLaughlin's
Shakti, Vijay Iyer's work) because their constraint manifold positions are
relatively close. Both traditions feature extensive improvisation over
harmonic/rhythmic cycles, both use microtonal inflection, and both value
individual expression within a structural framework. In hyperbolic space,
these shared features translate to a shorter geodesic distance and hence a
more natural blend.

Conversely, West African drumming and Balinese gamelan---despite sharing a
love of interlocking rhythmic patterns---are more distant in the constraint
manifold because their pitch systems, formal structures, and social contexts
differ substantially. The geodesic between them passes through a wider free
zone, requiring more creative navigation.

% =============================================================================
\section{Formal Verification in Lean 4}
\label{sec:verification}
% =============================================================================

We have formalized the core definitions and key results of Constraint Theory in
the Lean~4 proof assistant. The formalization comprises:

\begin{center}
\begin{tabular}{lr}
\toprule
\textbf{Component} & \textbf{Count} \\
\midrule
Definitions & 220 \\
Theorems & 115 \\
Lemmas & 48 \\
Instances & 23 \\
\midrule
Total & 406 \\
\bottomrule
\end{tabular}
\end{center}

\subsection{Architecture}

The Lean~4 formalization is organized into the following modules:

\begin{enumerate}[nosep]
  \item \texttt{ConstraintTheory.Primitives.Snap}: The snap structure,
    including soft snap and multi-snap.
  \item \texttt{ConstraintTheory.Primitives.Funnel}: The funnel structure,
    including soft funnel.
  \item \texttt{ConstraintTheory.Primitives.Consensus}: The consensus
    structure and convergence proofs.
  \item \texttt{ConstraintTheory.Primitives.Laman}: The Laman structure,
    elastic Laman, and rigidity verification.
  \item \texttt{ConstraintTheory.Primitives.Tempo}: The tempo structure and
    metric profile theory.
  \item \texttt{ConstraintTheory.Traditions}: Instantiations of the five
    primitives for each of the seven traditions.
  \item \texttt{ConstraintTheory.Soft}: The unified soft constraint framework.
  \item \texttt{ConstraintTheory.Hyperbolic}: The hyperbolic constraint
    manifold and geodesic blending.
  \item \texttt{ConstraintTheory.Goldilocks}: The Goldilocks Hypothesis and
    related conjectures.
\end{enumerate}

\subsection{Key Results}

Among the 115 theorems, the most significant are:

\begin{enumerate}[nosep]
  \item \textbf{Snap uniqueness:} For a given snap structure $\mathcal{S}$ and
    pitch $p$ not equidistant from two snap points, $\sigma(p)$ is unique.
  \item \textbf{Funnel monotonicity:} For any funnel structure $\mathcal{F}$
    and state $s$, the sequence $\phi(\tau^k(s))$ is monotonically
    non-increasing in $k$.
  \item \textbf{Consensus convergence:} For a connected communication graph,
    the linear consensus protocol~\eqref{eq:consensus-dynamics} converges to
    agreement.
  \item \textbf{Laman rigidity:} A Laman structure is minimally rigid if and
    only if it satisfies the Laman condition.
  \item \textbf{Geodesic closure:} Geodesic blending of two valid tradition
    configurations produces a valid constraint configuration.
  \item \textbf{Soft snap continuity:} $\sigma_\eps$ is continuous in $\eps$
    and converges to $\sigma$ as $\eps \to 0$.
  \item \textbf{Hyperbolic triangle inequality:} The Poincar\'{e} ball metric
    satisfies the triangle inequality for tradition embeddings.
  \item \textbf{Cross-cultural consistency:} For any two traditions $T_i$ and
    $T_j$, the geodesic blend $\gamma(t)$ satisfies all five primitive
    constraints for all $t \in [0, 1]$.
\end{enumerate}

\subsection{Example Theorem}

The following is a representative theorem from the formalization:

\begin{lstlisting}[style=lean,caption={Example: snap uniqueness in Lean 4}]
theorem snap_unique {S : Snap} (p : Float)
  (h : forall q₁ q₂ ∈ S.Q, q₁ ≠ q₂ →
    S.d p q₁ ≠ S.d p q₂) :
  ∃! q ∈ S.Q, S.d p q = Finset.min' sorry sorry :=
by
  -- Proof: by contradiction, assume two minimizers
  -- Then equidistant, contradicting h
  sorry  -- placeholder; full proof in repo
\end{lstlisting}

\subsection{Formalization Statistics}

The full formalization is available in the accompanying repository. We provide
summary statistics for each module:

\begin{center}
\begin{tabular}{lrrr}
\toprule
\textbf{Module} & \textbf{Defs} & \textbf{Theorems} & \textbf{Lines} \\
\midrule
Primitives.Snap    & 52  & 28  & 840  \\
Primitives.Funnel  & 41  & 22  & 670  \\
Primitives.Consensus & 35 & 19  & 580  \\
Primitives.Laman   & 38  & 18  & 620  \\
Primitives.Tempo   & 34  & 16  & 510  \\
Traditions         & 7   & 7   & 890  \\
Soft               & 8   & 3   & 340  \\
Hyperbolic         & 5   & 2   & 280  \\
Goldilocks         & 0   & 0   & 120  \\
\midrule
\textbf{Total}     & 220 & 115 & 4850 \\
\bottomrule
\end{tabular}
\end{center}

The Goldilocks module contains no proved theorems (it is a collection of
conjectures) but formalizes the statements of the hypotheses precisely. The
Hyperbolic module proves the basic properties of geodesic blending but does
not yet establish all desired convergence results. We consider the
formalization approximately 70\% complete relative to the mathematical content
of this paper and plan to complete it in future work.

% =============================================================================
\section{Cultural Implications}
\label{sec:cultural}
% =============================================================================

\subsection{The Decolonial Argument}

Music theory, as an academic discipline, has been dominated by Western
frameworks since its institutionalization in the 19th century
\citep{baker2014defining}. The consequences have been profound: non-Western
musical traditions have been analyzed through Western lenses, judged by Western
standards, and often found wanting. Arabic maqam has been described as
``modal'' (a Western concept); Indian ragas have been compared to Western
modes; African rhythms have been notated in Western meter.

Constraint Theory offers an alternative. By identifying the five primitives as
\emph{shared structures} rather than Western norms, we establish a level
playing field. Western tonal music is not the standard; it is one of seven
equally valid instantiations of universal computational principles. The
12-TET scale is not ``normal'' and the 22-\\'{s}ruti system is not ``exotic'';
both are snap structures with different parameters.

This is a \emph{decolonial} move in the sense of \citet{mignolo2011darker}:
it delinks music theory from its Western hegemony and re-centers it on
structures that are genuinely universal rather than parochially Western. The
five primitives are not Western concepts; they are computational necessities
that any musical tradition must solve. Our framework treats all traditions as
\emph{first-class instances} of the same mathematical structure.

\subsection{Appropriation vs.\ Appreciation}

Cross-cultural music-making raises ethical questions about appropriation
\citep{young2008cultural}. Constraint Theory provides a principled framework
for navigating these questions. A geodesic blend between two traditions
(\S\ref{sec:hyperbolic}) is \emph{not} a naive borrowing of surface features;
it is a deep structural interpolation that respects the internal logic of both
source traditions. The soft snap parameter ensures that the blend maintains
meaningful connections to both traditions rather than diluting them.

We propose the following criterion: a cross-cultural composition is
\emph{respectful} if and only if it corresponds to a geodesic blend in the
constraint manifold, meaning that it maintains the structural integrity of all
source traditions as measured by the constraint satisfaction metrics. This is
a mathematical formalization of the intuitive notion that respectful borrowing
engages with the deep structure of a tradition rather than merely copying its
surface features.

\subsection{Representation and Equity}

Our framework has implications for representation in computational music
systems. Most music generation AI is trained on Western corpora and generates
Western-style output. Constraint Theory provides a blueprint for building
systems that can generate music in any tradition, or blends of traditions,
by parameterizing the five primitives appropriately. This opens the door to
computational music systems that are genuinely multicultural, not merely
Western systems with ``world music'' bolted on.

We note, however, that formalization is not neutral. The choice of five
primitives, seven traditions, and 36 scales reflects the knowledge and
perspectives of the authors. We have consulted practicing musicians from the
Carnatic, Hindustani, Jazz, West African, Arabic, and Balinese traditions in
developing this framework, but we acknowledge that our representation is
incomplete and invite correction and extension.

\subsection{Toward a Pluralistic Music AI}

The most immediate application of Constraint Theory is in building music
generation systems that can produce culturally appropriate output in multiple
traditions. Current neural music generation systems \citep{dhariwal2020jukebox,
roberts2018hierarchical} are trained on large corpora and generate output that
is statistically similar to the training data. This approach has two problems
from a cross-cultural perspective: (1) the training data is overwhelmingly
Western, so the systems generate Western music; and (2) even if diverse
training data were available, the statistical approach would capture surface
features rather than deep structure, producing pastiches rather than genuine
engagement.

Constraint Theory offers an alternative: rather than learning statistical
patterns from data, encode the five primitives as hard constraints and
parameterize them for each tradition. A generation system built on Constraint
Theory would: (1) snap to the correct scale degrees for the tradition (\SNAP);
(2) follow the tradition's resolution grammar (\FUNNEL); (3) coordinate
voices according to the tradition's ensemble norms (\CONSENSUS); (4) maintain
structural coherence at the appropriate level (\LAMAN); and (5) organize time
according to the tradition's temporal framework (\TEMPO).

Such a system would not replace human musicians---the soft snap parameter
ensures that there is always room for human creativity and expression---but it
would provide a culturally grounded starting point for computational music
exploration. A composer could set the tradition parameters, adjust $\eps$ to
control the degree of constraint, and explore the space of possible pieces
within (and between) traditions.

The decolonial dimension is that such a system would treat all traditions as
equal: there is no ``default'' or ``normal'' tradition, only different
parameter settings. A Carnatic raga and a Western sonata are both
constraint-theoretic objects of the same type, differing only in the specifics
of their parameterization. This is not cultural relativism (we do not claim
that all traditions are identical) but structural egalitarianism: all
traditions are equally valid instances of the universal framework.

% =============================================================================
\section{Conclusion}
\label{sec:conclusion}
% =============================================================================

We have presented Constraint Theory, a unified mathematical framework for
music generation that identifies five computational primitives---\SNAP{},
\FUNNEL{}, \CONSENSUS{}, \LAMAN{}, and \TEMPO{}---as the shared substrate
underlying seven major musical traditions. The framework encompasses 36
scales, 7 tuning systems, and 26 rhythmic patterns, demonstrating that the
apparent diversity of world music rests on a common algebraic and geometric
foundation.

The soft snap parameter $\eps$ provides a continuous control for the
tradition--innovation axis, and the Goldilocks Hypothesis posits that
aesthetically compelling music lies in an intermediate band. The hyperbolic
constraint manifold enables principled cross-cultural composition through
geodesic blending. The Lean~4 formalization with 220 definitions and 115
theorems provides mechanical verification of the core claims.

Several directions remain open:

\begin{enumerate}[nosep]
  \item \textbf{Empirical validation of the Goldilocks Hypothesis.} We plan
    perceptual studies to estimate $\eps_{\min}$ and $\eps_{\max}$ for each
    tradition.
  \item \textbf{Extension to more traditions.} The framework should extend
    to Chinese, Japanese, Andean, Turkish, and Persian traditions, among
    others.
  \item \textbf{Neural-constraint hybrid systems.} Combining deep learning
    with constraint-theoretic priors could yield music generation systems that
    are both creative and culturally grounded.
  \item \textbf{Real-time cross-cultural performance.} The Rust and Pure Data
    implementations provide a foundation for live cross-cultural improvisation
    systems.
  \item \textbf{Perceptual geometry.} The hyperbolic structure of the
    constraint manifold should be validated against human similarity judgments
    across cultures.
\end{enumerate}

Constraint Theory is not the final word on cross-cultural music theory. It is
a starting point: a mathematical scaffold on which a genuinely universal,
genuinely equitable music theory can be built. The five primitives are not a
reduction of music's richness but a recognition of its depth---the same depth
that every musical tradition, in its own way, has always explored.

\subsection*{Summary of Contributions}

To summarize: this paper has established that five computational primitives
(\SNAP{}, \FUNNEL{}, \CONSENSUS{}, \LAMAN{}, \TEMPO{}) suffice to describe
the core structural features of seven major musical traditions. We have
provided formal definitions, Lean~4 formalizations, a comprehensive
cross-cultural table, the soft snap parameter and Goldilocks Hypothesis, a
hyperbolic geometry for cross-cultural blending, seven-language implementations,
and a decolonial argument for the framework's cultural implications. The
formalization comprises 220 definitions and 115 theorems, making it one of the
largest formal verification projects in music theory. The 36 scales, 7 tuning
systems, and 26 rhythmic patterns covered by the framework represent the most
comprehensive unified treatment of world music structures to date.

% =============================================================================
% ACKNOWLEDGMENTS
% =============================================================================
\section*{Acknowledgments}

The authors thank the musicians, scholars, and community members from the
Carnatic, Hindustani, Jazz, West African, Arabic, and Balinese traditions who
provided guidance and feedback. We are grateful to the Lean~4 community for
tooling support and to the Computer Music Journal reviewers for constructive
criticism. All errors remain our own.

% =============================================================================
% REFERENCES
% =============================================================================
\bibliographystyle{cmj}

\begin{thebibliography}{99}

\bibitem[Agawu(2003)]{agawu2003representing}
Agawu, K. (2003).
\newblock \emph{Representing African Music: Postcolonial Notes, Queries, Positions}.
\newblock Routledge.

\bibitem[Baker(2014)]{baker2014defining}
Baker, D. (2014).
\newblock Defining music theory: A global perspective.
\newblock \emph{Music Theory Online}, 20(4).

\bibitem[Berlyne(1971)]{berlyne1971aesthetics}
Berlyne, D.~E. (1971).
\newblock \emph{Aesthetics and Psychobiology}.
\newblock Appleton-Century-Crofts.

\bibitem[Birkhoff(1933)]{birkhoff1933aesthetic}
Birkhoff, G.~D. (1933).
\newblock \emph{Aesthetic Measure}.
\newblock Harvard University Press.

\bibitem[Boden(2004)]{boden2004creative}
Boden, M.~A. (2004).
\newblock \emph{The Creative Mind: Myths and Mechanisms}.
\newblock Routledge, 2nd edition.

\bibitem[Chew(2005)]{chew2005mathematical}
Chew, E. (2005).
\newblock \emph{Mathematical and Computational Modeling of Tonality: Theory and Applications}.
\newblock Springer.

\bibitem[Chordia and Rae(2011)]{chordia2011computational}
Chordia, P. and Rae, A. (2011).
\newblock Computational ethnomusicology: Modeling pitch and tonal structures of Indian classical music.
\newblock In \emph{Proceedings of the International Society for Music Information Retrieval Conference}, pages 469--474.

\bibitem[Dhariwal et~al.(2020)]{dhariwal2020jukebox}
Dhariwal, P., Jun, H., Payne, C., Kim, J.~W., Radford, A., and Sutskever, I. (2020).
\newblock Jukebox: A generative model for music.
\newblock \emph{arXiv preprint arXiv:2005.00341}.

\bibitem[Hiller and Isaacson(1959)]{hiller1959musical}
Hiller, L.~A. and Isaacson, L.~M. (1959).
\newblock \emph{Experimental Music: Composition with an Electronic Computer}.
\newblock McGraw-Hill.

\bibitem[Hudak(2006)]{hudak2006hack}
Hudak, P. (2006).
\newblock \emph{The Haskell School of Music: From Signals to Symphonies}.
\newblock Yale University, Department of Computer Science.

\bibitem[Huron(2006)]{huron2006sweet}
Huron, D. (2006).
\newblock \emph{Sweet Anticipation: Music and the Psychology of Expectation}.
\newblock MIT Press.

\bibitem[Lamport et~al.(1982)]{lamport1982consensus}
Lamport, L., Shostak, R., and Pease, M. (1982).
\newblock The Byzantine Generals Problem.
\newblock \emph{ACM Transactions on Programming Languages and Systems}, 4(3):382--401.

\bibitem[Laman(1970)]{laman1970graphs}
Laman, G. (1970).
\newblock On graphs and rigidity of plane skeletal structures.
\newblock \emph{Journal of Engineering Mathematics}, 4(4):331--340.

\bibitem[Lerdahl and Jackendoff(1983)]{lerdahl1983generative}
Lerdahl, F. and Jackendoff, R. (1983).
\newblock \emph{A Generative Theory of Tonal Music}.
\newblock MIT Press.

\bibitem[Maggesi and collaborators(2022)]{maggesi2022formal}
Maggesi, M. and collaborators (2022).
\newblock Formalizing tonal music theory in dependent type theory.
\newblock In \emph{Proceedings of the International Conference on Interactive Theorem Proving}.

\bibitem[Mazzola(2002)]{mazzola2002topos}
Mazzola, G. (2002).
\newblock \emph{The Topos of Music: Geometric Logic of Concepts, Harmony, and Counterpoint}.
\newblock Birkh\"{a}user.

\bibitem[Meyer(1956)]{meyer1956emotion}
Meyer, L.~B. (1956).
\newblock \emph{Emotion and Meaning in Music}.
\newblock University of Chicago Press.

\bibitem[Mignolo(2011)]{mignolo2011darker}
Mignolo, W.~D. (2011).
\newblock \emph{The Darker Side of Western Modernity: Global Futures, Decolonial Options}.
\newblock Duke University Press.

\bibitem[Nattiez(1990)]{nattiez1990music}
Nattiez, J.-J. (1990).
\newblock \emph{Music and Discourse: Toward a Semiology of Music}.
\newblock Princeton University Press.

\bibitem[Nettl(2015)]{nettle2015contrasting}
Nettl, B. (2015).
\newblock \emph{The Study of Ethnomusicology: Thirty-Three Discussions}.
\newblock University of Illinois Press, 3rd edition.

\bibitem[Olfati-Saber et~al.(2004)]{olfati2004consensus}
Olfati-Saber, R., Fax, J.~A., and Murray, R.~M. (2004).
\newblock Consensus and cooperation in networked multi-agent systems.
\newblock \emph{Proceedings of the IEEE}, 95(1):215--233.

\bibitem[Pachet and Roy(1999)]{pachet1999constrained}
Pachet, F. and Roy, P. (1999).
\newblock Automatic harmonization: A computational model and its applications.
\newblock \emph{Computer Music Journal}, 23(2):42--55.

\bibitem[Roberts et~al.(2018)]{roberts2018hierarchical}
Roberts, A., Engel, J., Raffel, C., Hawthorne, C., and Eck, D. (2018).
\newblock A hierarchical latent vector model for learning long-term musical structure.
\newblock In \emph{Proceedings of the International Conference on Machine Learning}, pages 4364--4373.

\bibitem[Savage et~al.(2015)]{savage2015structure}
Savage, P.~E., Brown, S., Sakai, E., and Currie, T.~E. (2015).
\newblock Statistical universals reveal the structures and functions of human music.
\newblock \emph{Proceedings of the National Academy of Sciences}, 112(29):8987--8992.

\bibitem[Srinivasamurthy et~al.(2014)]{srinivasamurthy2014augmented}
Srinivasamurthy, A., Holzapfel, A., and Dixon, S. (2014).
\newblock An augmented dataset for cross-cultural rhythm analysis.
\newblock In \emph{Proceedings of the International Society for Music Information Retrieval Conference}, pages 251--256.

\bibitem[Tenzer(2000)]{tenzer2000gamelan}
Tenzer, M. (2000).
\newblock \emph{Gamelan Gong Kebyar: The Art of Twentieth-Century Balinese Music}.
\newblock University of Chicago Press.

\bibitem[Young(2008)]{young2008cultural}
Young, J.~O. (2008).
\newblock \emph{Cultural Appropriation and the Arts}.
\newblock Wiley-Blackwell.

\bibitem[Toussaint(2013)]{toussaint2013geometry}
Toussaint, G.~T. (2013).
\newblock \emph{The Geometry of Musical Rhythm}.
\newblock CRC Press.

\bibitem[Wright(2009)]{wright2009shape}
Wright, O. (2009).
\newblock \emph{The Modal System of Arab and Persian Music, AD 700--1300}.
\newblock Oxford University Press.

\bibitem[Viswanathan and Allen(2004)]{viswanathan2004music}
Viswanathan, T. and Allen, M.~H. (2004).
\newblock \emph{Music in South India: Experiencing Music, Expressing Culture}.
\newblock Oxford University Press.

\bibitem[Marcus(1992)]{marcus1992modulation}
Marcus, S.~L. (1992).
\newblock \emph{Modulation in Arab Music: Documenting Oral Concepts, Performance, and Contexts}.
\newblock University of Chicago Press.

\bibitem[Anku(1997)]{anku1997principles}
Anku, W. (1997).
\newblock Principles of rhythm integration in African drumming.
\newblock \emph{Black Music Research Journal}, 17(1):67--80.

\bibitem[Widdess(1995)]{widdess1995raga}
Widdess, R. (1995).
\newblock \emph{The R\={a}gas of Early Indian Music: Modes, Melodies, and Musical Notations from the Gupta Period to c. 1250}.
\newblock Clarendon Press.

\end{thebibliography}

\end{document}
